\hypertarget{introduuxe7uxe3o}{%
\section{Introdução}\label{introduuxe7uxe3o}}

Este relatório apresenta a primeira entrega do seminário de testes de
características de qualidade da disciplina FGA0314, Testes de Software.
O tema escolhido pela equipe é \textbf{testes de autorização e controle
de acesso}, aplicado às APIs REST e à interface Web do projeto Educado
{[}1{]}.

Autorização é a subcaracterística de segurança que decide se um sujeito
já identificado pode executar uma operação sobre um recurso. Falhas
nesse controle ocupam a primeira posição do OWASP Top 10:2025 {[}2{]}, o
que justifica tratá-la como objeto de teste próprio, separado da
autenticação.

O documento está organizado da seguinte forma: a Seção 2 fundamenta os
conceitos de autorização e os modelos de controle de acesso; a Seção 3
apresenta as normas e referências usadas para avaliá-la; a Seção 4
descreve as técnicas de teste; a Seção 5 compara ferramentas de apoio; a
Seção 6 define métricas e uma proposta de aplicação ao Educado; a Seção
7 discute limitações e conclui.

Esta entrega é de pesquisa e planejamento. Nenhum teste foi executado no
Educado e nenhum resultado experimental é apresentado. Observações sobre
o código do projeto são tratadas como hipóteses de teste, não como
vulnerabilidades confirmadas.

\hypertarget{fundamentauxe7uxe3o-de-autorizauxe7uxe3o}{%
\section{Fundamentação de
autorização}\label{fundamentauxe7uxe3o-de-autorizauxe7uxe3o}}

\hypertarget{o-que-uxe9-autorizauxe7uxe3o}{%
\subsection{O que é autorização?}\label{o-que-uxe9-autorizauxe7uxe3o}}

Autorização é o processo de decidir se uma solicitação deve ser
permitida ou negada de acordo com uma política de acesso. Ela responde a
perguntas como ``este usuário pode consultar este prontuário?'' ou
``este serviço pode excluir este registro?''. O conceito não se confunde
com \textbf{identificação}, que associa uma entidade a um identificador,
nem com \textbf{autenticação}, que verifica se quem apresenta esse
identificador consegue demonstrar que é seu titular. Assim, informar um
nome de usuário identifica, apresentar uma credencial válida autentica e
avaliar o que a identidade autenticada pode fazer autoriza. A
autenticação pode fornecer dados confiáveis para a decisão, mas não
concede, por si só, acesso a todos os recursos {[}3{]}.

Uma decisão de acesso relaciona cinco elementos. O \textbf{sujeito} é a
entidade que solicita o acesso, como uma pessoa, processo ou serviço. O
\textbf{recurso} é o objeto protegido, por exemplo uma página, arquivo,
registro, campo ou função administrativa. A \textbf{operação} expressa a
ação pretendida, como ler, criar, alterar, aprovar ou excluir. A
\textbf{política} estabelece as condições sob as quais essa combinação é
permitida. Por fim, a \textbf{decisão de acesso} é o resultado da
avaliação --- normalmente permitir ou negar --- que deve ser
efetivamente aplicado pelo mecanismo de controle. Logo, autorizar não
significa apenas reconhecer um perfil geral: significa avaliar um
sujeito, uma operação e um recurso concretos no contexto da política
aplicável {[}3{]}.

\hypertarget{princuxedpios-de-proteuxe7uxe3o}{%
\subsection{Princípios de
proteção}\label{princuxedpios-de-proteuxe7uxe3o}}

Quatro princípios orientam um projeto de autorização seguro. O
\textbf{menor privilégio} determina que cada sujeito receba somente as
permissões necessárias para executar sua tarefa e apenas pelo tempo
necessário. Isso reduz tanto a superfície disponível a um invasor quanto
o dano causado por erro ou comprometimento de uma conta. A
\textbf{negação por padrão} exige que uma solicitação seja recusada
quando nenhuma regra a autorizar explicitamente; uma funcionalidade nova
ou uma combinação imprevista de condições não deve se tornar pública por
omissão {[}4{]}, {[}5{]}.

A \textbf{mediação completa} requer que toda tentativa de acesso seja
verificada, e não apenas a primeira ação de uma sessão. Uma tela
escondida, uma resposta anterior em cache ou a aprovação de acesso a um
recurso do mesmo tipo não substituem a verificação do objeto e da
operação atuais. Em aplicações Web, isso implica validar a autorização
no lado confiável do servidor em cada requisição, inclusive nas APIs,
sem depender de botões ocultos ou de verificações em JavaScript {[}4{]},
{[}5{]}.

A \textbf{separação de responsabilidades} evita concentrar poderes
incompatíveis em uma única pessoa ou função. Uma operação sensível pode
exigir condições independentes ou participação de sujeitos distintos;
por exemplo, quem cadastra um pagamento não deveria também aprová-lo.
Essa separação reduz fraude e abuso e pode ser expressa como restrições
de papéis ou como etapas diferentes de um fluxo. Ela complementa o menor
privilégio: além de limitar a quantidade de permissões, limita
combinações perigosas delas {[}4{]}, {[}6{]}.

\hypertarget{modelos-de-controle-de-acesso}{%
\subsection{Modelos de controle de
acesso}\label{modelos-de-controle-de-acesso}}

No \textbf{controle de acesso baseado em papéis} (\emph{Role-Based
Access Control}, RBAC), permissões são associadas a papéis que
representam funções organizacionais, como estudante, professor ou
administrador, e os sujeitos recebem um ou mais desses papéis. O modelo
simplifica a administração e favorece menor privilégio e separação de
responsabilidades, mas papéis excessivamente amplos podem continuar
concedendo acesso indevido {[}6{]}.

O \textbf{controle por propriedade} acrescenta uma relação entre sujeito
e recurso. Estar autenticado e poder usar a função ``editar'' não basta:
a aplicação confirma se o registro pertence ao solicitante ou está sob
sua responsabilidade. Por exemplo, um estudante pode alterar a própria
entrega, mas não a de outro estudante apenas trocando um identificador
na URL. A propriedade deve ser verificada no servidor e pode coexistir
com RBAC --- um professor, por exemplo, pode acessar entregas somente
das turmas pelas quais é responsável {[}2{]}.

No \textbf{controle de acesso baseado em atributos}
(\emph{Attribute-Based Access Control}, ABAC), a política avalia
atributos do sujeito, do recurso, da operação e, quando necessário, do
ambiente. Uma regra pode permitir a leitura quando o sujeito pertence ao
mesmo departamento do documento, possui credencial ativa e realiza a
solicitação durante determinado período. O ABAC representa condições
mais contextuais e granulares que papéis isolados, embora exija
atributos confiáveis e políticas cuidadosamente administradas {[}3{]}.

\hypertarget{por-que-uxe9-importante-em-aplicauxe7uxf5es-web}{%
\subsection{Por que é importante em aplicações
Web?}\label{por-que-uxe9-importante-em-aplicauxe7uxf5es-web}}

Aplicações Web expõem recursos e operações por rotas, formulários e APIs
acessíveis a clientes que não são confiáveis. O navegador mostra somente
a interface prevista, mas um solicitante pode alterar URLs,
identificadores, métodos HTTP e corpos de requisição. Por isso,
autenticar a sessão ou esconder controles visuais não protege os objetos
subjacentes. A autorização precisa acompanhar cada ponto de entrada e
aplicar as mesmas regras de negócio a páginas, APIs e recursos
estáticos. A OWASP mantém falhas de controle de acesso na primeira
posição de seu levantamento de riscos para aplicações Web, o que
evidencia a relevância prática do problema {[}2{]}.

\hypertarget{quais-problemas-podem-ocorrer}{%
\subsection{Quais problemas podem
ocorrer?}\label{quais-problemas-podem-ocorrer}}

Uma falha pode causar \textbf{exposição indevida}, quando alguém lê
dados pessoais, documentos ou informações de outra conta;
\textbf{alteração indevida}, quando modifica registros que não controla;
ou \textbf{exclusão indevida}, quando destrói dados por uma operação que
deveria estar restrita. Esses impactos comprometem, respectivamente,
confidencialidade, integridade e disponibilidade, além de poderem gerar
consequências legais e operacionais.

Também pode ocorrer \textbf{escalonamento horizontal de privilégios}, no
qual um usuário acessa recursos de outro usuário de nível semelhante, ou
\textbf{escalonamento vertical}, no qual obtém capacidades superiores,
como as de moderador ou administrador. No caso mais grave, a falha
permite executar operações administrativas: criar contas privilegiadas,
alterar políticas, aprovar transações, exportar bases ou remover
conteúdos. Referências diretas inseguras a objetos, endpoints sem
verificação e confiança em controles do cliente são caminhos comuns para
esses resultados {[}2{]}, {[}5{]}. Portanto, a autorização é a barreira
que transforma identidade autenticada em acesso limitado, justificável e
coerente com as responsabilidades de cada sujeito.

\hypertarget{normas-e-referuxeancias-para-avaliauxe7uxe3o-de-autorizauxe7uxe3o}{%
\section{Normas e referências para avaliação de
autorização}\label{normas-e-referuxeancias-para-avaliauxe7uxe3o-de-autorizauxe7uxe3o}}

\hypertarget{qualidade-do-produto-e-conscientizauxe7uxe3o}{%
\subsection{Qualidade do produto e
conscientização}\label{qualidade-do-produto-e-conscientizauxe7uxe3o}}

A ISO/\allowbreak{}IEC 25010:2023, segunda edição, integra a família SQuaRE e define
um modelo de qualidade de produto com nove características, subdivididas
em subcaracterísticas. O modelo orienta a especificação de requisitos,
os objetivos de teste e os critérios de aceitação durante o ciclo de
vida {[}7{]}. Para este seminário, a autorização é examinada no âmbito
da segurança: restringir a leitura protege a confidencialidade, enquanto
restringir alterações e exclusões protege a integridade. Essa relação
orienta a avaliação, mas precisa ser traduzida em regras concretas sobre
sujeitos, recursos e operações; a norma não constitui, por si só, um
roteiro de testes de endpoints.

O OWASP Top 10:2025 é uma referência de \textbf{conscientização} sobre
riscos de aplicações Web {[}8{]}. A categoria \textbf{A01 --- Broken
Access Control} reúne situações em que usuários ultrapassam suas
permissões, como acessar registros alheios, alcançar funções
administrativas ou contornar verificações manipulando requisições. Suas
recomendações incluem negação por padrão, verificação em código
confiável no servidor e controle da propriedade dos registros {[}9{]}. A
categoria ajuda a justificar prioridades, mas sua posição no ranking não
demonstra que o Educado tenha essas vulnerabilidades. Também não fornece
cobertura suficiente para declarar uma aplicação segura: \textbf{ASVS e
WSTG oferecem requisitos e procedimentos mais diretamente testáveis}.

\hypertarget{requisitos-selecionados-do-asvs}{%
\subsection{Requisitos selecionados do
ASVS}\label{requisitos-selecionados-do-asvs}}

O OWASP Application Security Verification Standard \textbf{5.0.0},
capítulo \textbf{V8 --- Authorization}, permite especificar o
comportamento esperado. Selecionam-se os requisitos abaixo; o prefixo da
versão evita ambiguidades entre edições {[}10{]}.

\begin{longtable}[]{@{}
  >{\raggedright\arraybackslash}p{(\columnwidth - 4\tabcolsep) * \real{0.3333}}
  >{\raggedright\arraybackslash}p{(\columnwidth - 4\tabcolsep) * \real{0.3333}}
  >{\raggedright\arraybackslash}p{(\columnwidth - 4\tabcolsep) * \real{0.3333}}@{}}
\toprule\noalign{}
\begin{minipage}[b]{\linewidth}\raggedright
Requisito
\end{minipage} & \begin{minipage}[b]{\linewidth}\raggedright
Nível
\end{minipage} & \begin{minipage}[b]{\linewidth}\raggedright
Verificação selecionada
\end{minipage} \\
\midrule\noalign{}
\endhead
\bottomrule\noalign{}
\endlastfoot
\texttt{v5.0.0-8.1.1} & 1 & Documentar regras de acesso a funções e
dados, considerando permissões e atributos dos recursos. \\
\texttt{v5.0.0-8.1.2} & 2 & Documentar restrições de leitura e escrita
por campo. \\
\texttt{v5.0.0-8.2.1} & 1 & Restringir funções a sujeitos com permissão
explícita. \\
\texttt{v5.0.0-8.2.2} & 1 & Restringir cada objeto aos sujeitos
autorizados. \\
\texttt{v5.0.0-8.2.3} & 2 & Restringir campos específicos conforme as
permissões. \\
\texttt{v5.0.0-8.3.1} & 1 & Aplicar as regras em camada confiável, sem
depender do cliente. \\
\end{longtable}

Os itens de documentação permitem definir o oráculo dos testes: uma
matriz de sujeitos, operações, objetos e campos com decisões esperadas.
Os demais orientam verificações positivas e negativas contra essa
matriz. A seleção combina requisitos dos níveis 1 e 2; não equivale à
verificação integral de qualquer nível do ASVS.

\hypertarget{procedimentos-do-wstg}{%
\subsection{Procedimentos do WSTG}\label{procedimentos-do-wstg}}

O OWASP Web Security Testing Guide \textbf{4.2}, seção 4.5, descreve
quatro procedimentos de autorização. Neste mapeamento, os
identificadores incluem a versão \texttt{v42} para preservar a
rastreabilidade.

\begin{itemize}
\tightlist
\item
  \textbf{WSTG-v42-ATHZ-01 --- Directory Traversal/\allowbreak{}File Include:}
  investigar parâmetros que referenciam arquivos e tentativas de acesso
  fora dos limites permitidos; aplica-se quando existirem
  funcionalidades de arquivos ou inclusão {[}11{]}.
\item
  \textbf{WSTG-v42-ATHZ-02 --- Bypassing Authorization Schema:} repetir
  requisições sem autenticação, após logout e com outras identidades ou
  papéis, procurando desvios horizontais e verticais; relaciona-se aos
  requisitos \texttt{8.2.1}, \texttt{8.2.2} e \texttt{8.3.1} {[}12{]}.
\item
  \textbf{WSTG-v42-ATHZ-03 --- Privilege Escalation:} manipular
  parâmetros de grupo ou perfil e verificar obtenção de capacidades
  indevidas; relaciona-se a \texttt{8.2.1} e, quando a alteração envolve
  campos protegidos, a \texttt{8.2.3} {[}13{]}.
\item
  \textbf{WSTG-v42-ATHZ-04 --- Insecure Direct Object References:}
  substituir referências de objetos usando contas com permissões
  conhecidas e verificar acesso indevido; operacionaliza especialmente
  \texttt{8.2.2} {[}14{]}.
\end{itemize}

Essas relações são uma proposta de organização do seminário, não uma
equivalência oficial ou exaustiva entre os documentos. No planejamento,
cada caso deve registrar precondições, identidade, recurso, operação e
resultado esperado. A avaliação deve considerar os dados retornados e os
efeitos persistidos, além do código HTTP: uma mensagem de erro não basta
se a operação proibida ocorrer.

\hypertarget{apis-e-classificauxe7uxe3o-das-fraquezas}{%
\subsection{APIs e classificação das
fraquezas}\label{apis-e-classificauxe7uxe3o-das-fraquezas}}

O \textbf{OWASP API Security Top 10:2023} detalha três dimensões
complementares. \textbf{API1 --- Broken Object Level Authorization}
trata da permissão sobre o objeto solicitado: poder chamar uma função
não autoriza consultar qualquer registro {[}15{]}. \textbf{API3 ---
Broken Object Property Level Authorization} trata das propriedades desse
objeto, tanto na exposição de campos quanto na alteração indevida;
corresponde ao foco de \texttt{8.2.3} {[}16{]}. \textbf{API5 --- Broken
Function Level Authorization} trata do acesso à própria função, como
executar uma operação administrativa com um perfil comum; corresponde ao
foco de \texttt{8.2.1} {[}17{]}. Assim, testar somente a troca de
identificadores deixa descobertas as permissões sobre campos e funções.

Como vocabulário de classificação, adota-se o catálogo \textbf{CWE 5.0},
preservado em arquivo versionado {[}18{]}. \textbf{CWE-862} identifica
ausência de verificação de autorização; \textbf{CWE-863}, uma
verificação existente, porém incorreta; \textbf{CWE-639}, acesso a dados
de outro usuário pela alteração de uma chave controlável pelo
solicitante {[}19{]}, {[}20{]}, {[}21{]}. O último se relaciona a
cenários de API1 e ATHZ-04, sem tornar BOLA, IDOR e CWE-639 sinônimos
universais. A classificação deve acompanhar a causa demonstrada,
evitando atribuir automaticamente todas as categorias a um mesmo caso.

\hypertarget{comparauxe7uxe3o-e-uso-conjunto}{%
\subsection{Comparação e uso
conjunto}\label{comparauxe7uxe3o-e-uso-conjunto}}

A tabela sintetiza o papel das referências apresentadas.

\begin{longtable}[]{@{}
  >{\raggedright\arraybackslash}p{(\columnwidth - 6\tabcolsep) * \real{0.2500}}
  >{\raggedright\arraybackslash}p{(\columnwidth - 6\tabcolsep) * \real{0.2500}}
  >{\raggedright\arraybackslash}p{(\columnwidth - 6\tabcolsep) * \real{0.2500}}
  >{\raggedright\arraybackslash}p{(\columnwidth - 6\tabcolsep) * \real{0.2500}}@{}}
\toprule\noalign{}
\begin{minipage}[b]{\linewidth}\raggedright
Referência e versão
\end{minipage} & \begin{minipage}[b]{\linewidth}\raggedright
Objetivo
\end{minipage} & \begin{minipage}[b]{\linewidth}\raggedright
Escopo
\end{minipage} & \begin{minipage}[b]{\linewidth}\raggedright
Uso na avaliação
\end{minipage} \\
\midrule\noalign{}
\endhead
\bottomrule\noalign{}
\endlastfoot
ISO/\allowbreak{}IEC 25010:2023 & Organizar a qualidade & Produtos de software e TIC
& Definir objetivos e critérios de qualidade. \\
OWASP Top 10:2025 & Conscientizar sobre riscos & Aplicações Web; A01 &
Justificar prioridade para autorização. \\
ASVS 5.0.0, V8 & Especificar requisitos verificáveis & Documentação e
controles de autorização & Definir resultados esperados e cobertura. \\
WSTG 4.2, seção 4.5 & Orientar procedimentos & Testes de segurança Web &
Estruturar tentativas de contorno. \\
API Security Top 10:2023 & Evidenciar riscos de APIs & Objetos,
propriedades e funções & Diversificar cenários de autorização. \\
CWE 5.0: 862, 863 e 639 & Classificar fraquezas & Causas e mecanismos de
falha & Padronizar a descrição de eventuais achados. \\
\end{longtable}

Para a futura aplicação ao Educado, propõe-se encadear objetivo de
qualidade, risco, requisito, procedimento e classificação. A política
real do projeto deverá determinar quais combinações são permitidas e
quais casos se aplicam. Esta entrega estabelece referências para esse
planejamento, sem executar testes nem afirmar resultados experimentais.

\hypertarget{tuxe9cnicas-de-teste-de-autorizauxe7uxe3o-e-controle-de-acesso}{%
\section{Técnicas de teste de autorização e controle de
acesso}\label{tuxe9cnicas-de-teste-de-autorizauxe7uxe3o-e-controle-de-acesso}}

\hypertarget{modelo-de-teste-e-matriz-de-decisuxe3o}{%
\subsection{Modelo de teste e matriz de
decisão}\label{modelo-de-teste-e-matriz-de-decisuxe3o}}

O teste começa por uma matriz com \textbf{sujeito}, \textbf{recurso},
\textbf{operação} e \textbf{resultado esperado}. Por exemplo, para um
estudante (\texttt{S1}), um professor (\texttt{S2}) e um administrador
(\texttt{S3}):

\begin{longtable}[]{@{}
  >{\raggedright\arraybackslash}p{(\columnwidth - 6\tabcolsep) * \real{0.2500}}
  >{\raggedright\arraybackslash}p{(\columnwidth - 6\tabcolsep) * \real{0.2500}}
  >{\raggedright\arraybackslash}p{(\columnwidth - 6\tabcolsep) * \real{0.2500}}
  >{\raggedright\arraybackslash}p{(\columnwidth - 6\tabcolsep) * \real{0.2500}}@{}}
\toprule\noalign{}
\begin{minipage}[b]{\linewidth}\raggedright
Sujeito
\end{minipage} & \begin{minipage}[b]{\linewidth}\raggedright
Recurso
\end{minipage} & \begin{minipage}[b]{\linewidth}\raggedright
Operação
\end{minipage} & \begin{minipage}[b]{\linewidth}\raggedright
Resultado esperado
\end{minipage} \\
\midrule\noalign{}
\endhead
\bottomrule\noalign{}
\endlastfoot
Anônimo & Perfil de estudante & Ler & Negar (401 ou redirecionamento
para login) \\
\texttt{S1} & Sua própria entrega & Ler/\allowbreak{}alterar & Permitir \\
\texttt{S1} & Entrega de outro estudante & Ler/\allowbreak{}alterar & Negar (403 ou
resposta indistinguível) \\
\texttt{S2} & Entrega de estudante de sua turma & Ler/\allowbreak{}avaliar &
Permitir \\
\texttt{S2} & Função de administração & Excluir usuário & Negar \\
\texttt{S3} & Usuário e política & Administrar & Permitir, com controles
adicionais para operações críticas \\
\end{longtable}

A matriz diferencia permissão funcional de propriedade. O caso registra
precondições, usuário, token ou sessão, recurso, requisição, resultado
esperado e evidência (status, corpo, efeito persistido e logs). Um
\texttt{403} isolado não basta: é necessário verificar que o dado não
foi retornado nem alterado. Toda execução deve usar dados descartáveis
em ambiente local ou de homologação {[}5{]}, {[}22{]}.

\hypertarget{acesso-anuxf4nimo-e-navegauxe7uxe3o-foruxe7ada}{%
\subsection{Acesso anônimo e navegação
forçada}\label{acesso-anuxf4nimo-e-navegauxe7uxe3o-foruxe7ada}}

\textbf{Objetivo:} confirmar que recursos protegidos exigem autenticação
e autorização no servidor. \textbf{Procedimento:} sem cookies, tokens ou
cabeçalhos de autenticação, acessar diretamente páginas, APIs e arquivos
previsíveis; depois repetir com conta autenticada, mas sem o papel
necessário. Testar a rota mesmo quando o menu não a apresenta.
\textbf{Falha detectada:} conteúdo privado entregue a anônimos, endpoint
que aceita uma sessão insuficiente ou confiança no redirecionamento e no
controle visual do cliente. \textbf{Vantagem:} é simples, barato e
encontra rotas esquecidas. \textbf{Limitação:} não cobre relações entre
objetos ou regras condicionais.

O frontend pode esconder um botão, mas o usuário consegue construir a
requisição fora do navegador. A decisão confiável é a aplicada pelo
servidor em cada requisição {[}5{]}.

\hypertarget{escalonamento-horizontal-bolaidor-e-adulterauxe7uxe3o-de-identificadores}{%
\subsection{Escalonamento horizontal, BOLA/\allowbreak{}IDOR e adulteração de
identificadores}\label{escalonamento-horizontal-bolaidor-e-adulterauxe7uxe3o-de-identificadores}}

\textbf{Objetivo:} verificar se um usuário com a mesma função acessa
apenas os próprios objetos. \textbf{Procedimento:} criar dois usuários
do mesmo nível, obter um recurso de cada um e substituir, em URL, query,
corpo ou caminho, o identificador pelo do outro. Repetir para leitura,
atualização e exclusão, usando IDs sequenciais e UUIDs. \textbf{Falha
detectada:} BOLA/\allowbreak{}IDOR, quando o servidor aceita a referência ao objeto
sem conferir a relação entre sujeito e recurso. \textbf{Vantagem:}
reproduz falha de alto impacto com alteração pequena e observável.
\textbf{Limitação:} exige dados de usuários distintos e pode não revelar
autorização contextual.

Testar também IDs em exportações e anexos; identificador difícil não é
autorização. O resultado para o objeto alheio é uma negativa
consistente, sem vazamento de existência, metadados ou efeitos {[}2{]},
{[}23{]}.

\hypertarget{escalonamento-vertical-e-fluxos-administrativos}{%
\subsection{Escalonamento vertical e fluxos
administrativos}\label{escalonamento-vertical-e-fluxos-administrativos}}

\textbf{Objetivo:} impedir que um usuário comum execute operações
reservadas a papéis superiores. \textbf{Procedimento:} comparar as
matrizes de estudante, professor e administrador; repetir requisições
administrativas com o token de menor privilégio; tentar acessar rotas de
gestão, alterar papéis e excluir contas. Para cada negativa, verificar
que não houve mudança no banco ou em fila. \textbf{Falha detectada:}
escalonamento vertical, endpoint administrativo sem verificação de
papel, papel amplo demais ou separação de responsabilidades ausente.
\textbf{Vantagem:} testa o impacto mais severo. \textbf{Limitação:}
aprovações múltiplas e expiração exigem cenários longos.

Testar a alteração do próprio papel e a criação de outro administrador,
além da operação final. Alterar o papel enviado ou guardado no cliente
(campo \texttt{role}, \texttt{localStorage} ou token não assinado) e
confirmar que o servidor o ignora. Ações críticas podem exigir
reautenticação ou segunda autoridade {[}5{]}.

\hypertarget{muxe9todo-http-e-autorizauxe7uxe3o-por-campo}{%
\subsection{Método HTTP e autorização por
campo}\label{muxe9todo-http-e-autorizauxe7uxe3o-por-campo}}

\textbf{Objetivo:} verificar que a autorização permanece correta quando
a forma da requisição muda. \textbf{Procedimento:} para recurso somente
de leitura, trocar \texttt{GET} por \texttt{POST}, \texttt{PUT},
\texttt{PATCH} ou \texttt{DELETE}; testar
\texttt{X-HTTP-Method-Override} e variantes de rota. Adicionar ao JSON
campos como \texttt{role}, \texttt{ownerId}, \texttt{approved} ou
\texttt{status} e tentar alterá-los com operação permitida.
\textbf{Falha detectada:} autorização aplicada a uma rota ou método, mas
não a outro, e \emph{mass assignment} que permite ao cliente controlar
atributos reservados. \textbf{Vantagem:} encontra inconsistências entre
controladores e mapeamentos. \textbf{Limitação:} requer conhecer o
contrato.

O resultado esperado é validar cada combinação de operação e campo no
servidor, aceitando apenas atributos permitidos para aquele sujeito. Não
se deve confundir filtragem visual do formulário com uma lista de campos
autorizados {[}24{]}.

\hypertarget{testes-manuais-automatizados-e-orientados-por-openapi}{%
\subsection{Testes manuais, automatizados e orientados por
OpenAPI}\label{testes-manuais-automatizados-e-orientados-por-openapi}}

Testes manuais exploram fluxos novos, interpretam regras de negócio e
confirmam evidências. Testes automatizados repetem a matriz para
usuários e identificadores, detectam regressões e comparam status, corpo
e efeitos. Testes orientados por OpenAPI geram casos de rotas, métodos,
parâmetros e esquemas documentados. Sua limitação é estrutural: um
contrato incompleto não descreve autorização, propriedade ou transições
de negócio, e a geração não substitui cenários com contas relacionadas
{[}25{]}.

Uma estratégia equilibrada automatiza casos de alto risco, usa a
especificação como inventário e reserva exploração manual para regras
contextuais. Cada falha deve preservar requisição e resposta
sanitizadas, usuário de teste e estado inicial, sem expor segredos.
Ferramentas devem ficar em ambientes locais ou de homologação, com
autorização explícita e taxa controlada. O objetivo é demonstrar a
decisão de acesso e seu impacto, não contornar controles em sistemas de
terceiros.

\hypertarget{ferramentas-para-testes-de-autorizauxe7uxe3o}{%
\section{Ferramentas para testes de
autorização}\label{ferramentas-para-testes-de-autorizauxe7uxe3o}}

Testar autorização exige observar a aplicação como cliente externo, por
requisições em uma suíte ou pelo comportamento da interface. Este texto
compara cinco alternativas --- OWASP ZAP com Access Control Testing,
Playwright com \texttt{APIRequestContext}, Vitest, Schemathesis e Burp
Suite Community com Autorize --- quanto a funcionamento, licença,
automação, CI, múltiplos usuários, vantagens e limitações, priorizando
opções open source alinhadas ao Educado {[}1{]}.

\hypertarget{owasp-zap-e-o-add-on-access-control-testing}{%
\subsection{OWASP ZAP e o add-on Access Control
Testing}\label{owasp-zap-e-o-add-on-access-control-testing}}

O OWASP ZAP é um proxy DAST sob Apache 2.0 {[}26{]}. O add-on Access
Control Testing define usuários e autenticação, cadastra regras por
contexto e tenta acessar cada URL como cada usuário, sinalizando
violações {[}27{]}. A automação usa Automation Framework, API, CLI e
Docker; o CI é viável com GitHub Action {[}26{]}. Múltiplos usuários são
nativos. As vantagens são cobertura ampla e detecção automatizada. As
limitações incluem configuração de contexto, dependência de exploração
prévia em SPAs e falsos positivos ou negativos {[}27{]}.

\hypertarget{playwright-com-apirequestcontext}{%
\subsection{Playwright com
APIRequestContext}\label{playwright-com-apirequestcontext}}

O Playwright é uma biblioteca de automação de navegadores e testes de
API sob Apache 2.0 {[}28{]}. O \texttt{APIRequestContext} dispara HTTP
sem abrir página, em contextos isolados com cookies e cabeçalhos
próprios {[}29{]}, permitindo comparar perfis: um estudante recebe
\texttt{403} ao alterar recurso alheio; um professor autorizado recebe
\texttt{200}. A automação ocorre por CLI, com CI por ações oficiais
{[}28{]}. Múltiplos usuários decorrem de contextos isolados e do estado
de autenticação. As vantagens são API e interface na mesma ferramenta e
controle por requisição. As limitações são a ausência de descoberta
automática e a manutenção manual.

\hypertarget{vitest-para-guards-e-comportamentos-do-frontend}{%
\subsection{Vitest para guards e comportamentos do
frontend}\label{vitest-para-guards-e-comportamentos-do-frontend}}

O Vitest é um executor de testes unitários e de componentes sobre Vite,
sob licença MIT {[}30{]}. Roda em jsdom e simula funções, módulos e
requisições, sendo adequado para guards de rota, renderização por perfil
e mensagens de negação {[}30{]}. A automação ocorre por
\texttt{npm\ test} ou \texttt{vitest\ run}, e o CI executa o script.
Múltiplos usuários são simulados por mocks. As vantagens são velocidade,
integração com o \texttt{educado-web} e baixo custo de regressão. As
limitações são decisivas: não prova autorização no servidor, mocks podem
divergir da API e a ferramenta não substitui um teste de segurança.

\hypertarget{schemathesis-para-operauxe7uxf5es-descritas-por-openapi}{%
\subsection{Schemathesis para operações descritas por
OpenAPI}\label{schemathesis-para-operauxe7uxf5es-descritas-por-openapi}}

O Schemathesis gera testes baseados em propriedades a partir de OpenAPI,
GraphQL ou gRPC, sob licença MIT {[}31{]}. Ele produz entradas válidas e
inválidas, verifica erros de servidor e violações de contrato e, com
verificações personalizadas, compara respostas entre credenciais para
detectar autorização quebrada em nível de objeto {[}31{]}. A automação
ocorre por CLI ou \texttt{pytest}, com CI por ação oficial {[}31{]}.
Múltiplos usuários dependem de autenticação configurada e de
verificações próprias. As vantagens são geração automática e achados
reproduzíveis. As limitações são a exigência de especificação completa,
verificações manuais e efeitos colaterais nos dados.

\hypertarget{burp-suite-community-e-autorize}{%
\subsection{Burp Suite Community e
Autorize}\label{burp-suite-community-e-autorize}}

O Burp Suite Community é proprietário e gratuito, voltado ao teste
manual, sem o scanner da versão Professional {[}32{]}. O Autorize é uma
extensão open source que repete requisições com baixo privilégio ou sem
autenticação e compara respostas, apoiando a detecção de IDOR {[}33{]}.
A automação limita-se à interceptação, sem CI oficial, mas há suporte a
vários perfis. Vantagens: uso exploratório; limitações: heurísticas,
ajuste manual e ausência de pipeline.

\hypertarget{comparauxe7uxe3o}{%
\subsection{Comparação}\label{comparauxe7uxe3o}}

\begin{longtable}[]{@{}
  >{\raggedright\arraybackslash}p{(\columnwidth - 8\tabcolsep) * \real{0.2000}}
  >{\raggedright\arraybackslash}p{(\columnwidth - 8\tabcolsep) * \real{0.2000}}
  >{\raggedright\arraybackslash}p{(\columnwidth - 8\tabcolsep) * \real{0.2000}}
  >{\raggedright\arraybackslash}p{(\columnwidth - 8\tabcolsep) * \real{0.2000}}
  >{\raggedright\arraybackslash}p{(\columnwidth - 8\tabcolsep) * \real{0.2000}}@{}}
\toprule\noalign{}
\begin{minipage}[b]{\linewidth}\raggedright
Ferramenta
\end{minipage} & \begin{minipage}[b]{\linewidth}\raggedright
Licença
\end{minipage} & \begin{minipage}[b]{\linewidth}\raggedright
Automação/\allowbreak{}CI
\end{minipage} & \begin{minipage}[b]{\linewidth}\raggedright
Múltiplos usuários
\end{minipage} & \begin{minipage}[b]{\linewidth}\raggedright
Limitação
\end{minipage} \\
\midrule\noalign{}
\endhead
\bottomrule\noalign{}
\endlastfoot
ZAP + Access Control Testing & Apache 2.0 & CLI, API, Docker e GitHub
Action & Sim, nativo & Configuração e falsos positivos \\
Playwright + APIRequestContext & Apache 2.0 & CLI e ações de CI & Sim,
contextos isolados & Não descobre falhas \\
Vitest & MIT & \texttt{vitest\ run} no pipeline & Simulado por mocks &
Não valida o servidor \\
Schemathesis & MIT & CLI, \texttt{pytest} e ação de CI & Configurável &
Depende da especificação \\
Burp Community + Autorize & Proprietária e open source & Manual; sem CI
oficial & Sim, vários perfis & Heurístico \\
\end{longtable}

\hypertarget{recomendauxe7uxe3o}{%
\subsection{Recomendação}\label{recomendauxe7uxe3o}}

A estratégia principal combina Playwright e OWASP ZAP. O Playwright
cobre cenários conhecidos e repetíveis, comparando respostas entre
perfis na API e na interface {[}29{]}. O ZAP percorre a aplicação e
confronta o acesso de cada usuário com as regras declaradas, encontrando
rotas esquecidas {[}27{]}. O Vitest complementa o frontend com guards e
mensagens de negação, sem provar a segurança do backend {[}30{]}.
Schemathesis é útil quando houver OpenAPI confiável {[}31{]}, e
Burp/\allowbreak{}Autorize fica como alternativa manual não prioritária {[}33{]}.
Nenhuma ferramenta foi instalada ou executada.

\hypertarget{muxe9tricas-e-proposta-de-aplicauxe7uxe3o}{%
\section{Métricas e Proposta de
aplicação}\label{muxe9tricas-e-proposta-de-aplicauxe7uxe3o}}

\hypertarget{objetivo-e-fronteira}{%
\subsection{Objetivo e fronteira}\label{objetivo-e-fronteira}}

Avaliar se a API REST do Educado permite que cada ator execute somente
as operações e acesse somente os dados autorizados para seu papel e sua
relação com cada objeto. O foco é autorização horizontal (um usuário
acessando dados de outro), vertical (um papel de menor privilégio
executando funções administrativas), adulteração de identificadores e de
campos sensíveis.

A autorização efetiva deve ser validada na
\textbf{\texttt{educado-api}}, diretamente nos endpoints HTTP.

\hypertarget{atores-de-teste}{%
\subsubsection{Atores de teste}\label{atores-de-teste}}

\begin{longtable}[]{@{}
  >{\raggedright\arraybackslash}p{(\columnwidth - 6\tabcolsep) * \real{0.2500}}
  >{\raggedright\arraybackslash}p{(\columnwidth - 6\tabcolsep) * \real{0.2500}}
  >{\raggedright\arraybackslash}p{(\columnwidth - 6\tabcolsep) * \real{0.2500}}
  >{\raggedright\arraybackslash}p{(\columnwidth - 6\tabcolsep) * \real{0.2500}}@{}}
\toprule\noalign{}
\begin{minipage}[b]{\linewidth}\raggedright
Ator
\end{minipage} & \begin{minipage}[b]{\linewidth}\raggedright
Identificação de teste
\end{minipage} & \begin{minipage}[b]{\linewidth}\raggedright
Papel/\allowbreak{}token
\end{minipage} & \begin{minipage}[b]{\linewidth}\raggedright
Uso
\end{minipage} \\
\midrule\noalign{}
\endhead
\bottomrule\noalign{}
\endlastfoot
Anônimo & -- & Sem token & Verificar recursos protegidos e recursos
públicos limitados \\
\texttt{USER} & U1, criador de conteúdo & \texttt{USER} & Acessar perfil
e conteúdo que possui; tentar ações de outro proprietário e
administrativas \\
\texttt{ADMIN} & ADM1 & \texttt{ADMIN} & Verificar funções
administrativas e acesso transversal explicitamente autorizado \\
Estudante A & A & \texttt{STUDENT} & Consultar/\allowbreak{}alterar somente perfil,
matrícula, progresso e certificados próprios \\
Estudante B & B & \texttt{STUDENT} & Identidade diferente de A com o
mesmo papel; fornecer objetos de B para testes de isolamento \\
\end{longtable}

\hypertarget{poluxedtica-inicial-e-matriz-de-autorizauxe7uxe3o}{%
\subsection{Política inicial e matriz de
autorização}\label{poluxedtica-inicial-e-matriz-de-autorizauxe7uxe3o}}

\hypertarget{legenda}{%
\subsubsection{Legenda}\label{legenda}}

\begin{itemize}
\tightlist
\item
  \textbf{Público limitado:} leitura de catálogo ou verificação pública,
  sem dados privados nem operações de escrita.
\item
  \textbf{Próprio:} somente o próprio perfil/\allowbreak{}objeto do ator autenticado.
\item
  \textbf{Próprio curso:} autor \texttt{USER} pode atuar nos cursos,
  seções, atividades e mídias que possui; \texttt{ADMIN} pode atuar em
  qualquer um.
\item
  \textbf{Elegível:} estudante pode consultar conteúdo disponibilizado
  para estudantes e responder atividades permitidas.
\item
  \textbf{Admin:} somente funções administrativas explicitamente
  previstas.
\item
  \textbf{Negado:} a operação não deve ser autorizada.
\end{itemize}

\begin{longtable}[]{@{}
  >{\raggedright\arraybackslash}p{(\columnwidth - 10\tabcolsep) * \real{0.1667}}
  >{\raggedright\arraybackslash}p{(\columnwidth - 10\tabcolsep) * \real{0.1667}}
  >{\raggedright\arraybackslash}p{(\columnwidth - 10\tabcolsep) * \real{0.1667}}
  >{\raggedright\arraybackslash}p{(\columnwidth - 10\tabcolsep) * \real{0.1667}}
  >{\raggedright\arraybackslash}p{(\columnwidth - 10\tabcolsep) * \real{0.1667}}
  >{\raggedright\arraybackslash}p{(\columnwidth - 10\tabcolsep) * \real{0.1667}}@{}}
\toprule\noalign{}
\begin{minipage}[b]{\linewidth}\raggedright
Recurso
\end{minipage} & \begin{minipage}[b]{\linewidth}\raggedright
Anônimo
\end{minipage} & \begin{minipage}[b]{\linewidth}\raggedright
\texttt{USER} U1
\end{minipage} & \begin{minipage}[b]{\linewidth}\raggedright
\texttt{ADMIN}
\end{minipage} & \begin{minipage}[b]{\linewidth}\raggedright
Estudante A
\end{minipage} & \begin{minipage}[b]{\linewidth}\raggedright
Estudante B
\end{minipage} \\
\midrule\noalign{}
\endhead
\bottomrule\noalign{}
\endlastfoot
Perfil & Negado & Ler/\allowbreak{}alterar o próprio perfil em \texttt{/me/\allowbreak{}profile} &
Ler/\allowbreak{}alterar o próprio perfil; consultar usuários apenas por funções
admin autorizadas & Ler/\allowbreak{}alterar apenas o perfil A & Ler/\allowbreak{}alterar apenas o
perfil B \\
Cursos & Público limitado pelo catálogo (\texttt{/catalog}); sem rotas
de autoria & Ler/\allowbreak{}criar/\allowbreak{}alterar/\allowbreak{}ativar/\allowbreak{}desativar/\allowbreak{}excluir somente cursos
próprios & Operar cursos de qualquer autor nas funções admin previstas &
Ler catálogo e conteúdo elegível; sem autoria & Igual a A, sem acesso ao
estado privado de A \\
Atividades & Leitura pública de atividades publicadas no catálogo, sem
expor o gabarito antes da resposta; escrita negada &
Criar/\allowbreak{}ler/\allowbreak{}alterar/\allowbreak{}excluir somente em curso próprio &
Criar/\allowbreak{}ler/\allowbreak{}alterar/\allowbreak{}excluir em qualquer curso & Consumir e responder
atividades elegíveis; sem CRUD nem acesso ao gabarito antes da resposta
& Igual a A; sem alterar respostas/\allowbreak{}progresso de A \\
Mídias & Fazer stream de mídia ativa e publicada; sem acesso a metadados
privados & Gerenciar metadados e excluir mídia própria; fazer stream de
qualquer mídia ativa e publicada & Administrar mídias conforme endpoint
admin; fazer stream de mídia ativa e publicada & Fazer stream de mídia
ativa e publicada; sem alterar/\allowbreak{}excluir & Igual a A; não acessar
metadados privados de outros autores \\
Instituições & Negado & Negado & Listar, consultar, criar, alterar e
excluir & Negado & Negado \\
Usuários & Negado; exceção para cadastro/\allowbreak{}login/\allowbreak{}verificação pública
previstos & Consultar/\allowbreak{}alterar somente os campos próprios; não
administrar papéis ou outros usuários & Listar/\allowbreak{}consultar e executar
ações admin previstas, como alterar papel ou excluir usuário & Somente
dados próprios permitidos & Somente dados próprios permitidos \\
Matrículas & Negado & Negado & \texttt{ADMIN}: acesso administrativo a
matrículas; discutir com a equipe quais operações e dados são permitidos
& Listar, criar, consultar e cancelar somente matrículas de A & Somente
matrículas de B \\
Progresso & Negado & Negado & \texttt{ADMIN}: acesso administrativo ao
progresso; discutir com a equipe o escopo e as operações permitidas &
Consultar/\allowbreak{}gravar/\allowbreak{}concluir somente progresso de A & Somente progresso de
B \\
Certificados & No código atual, sem autenticação:
\texttt{GET\ /certificates/\{username\}} lista certificados,
\texttt{POST\ /certificates} cria certificado e
\texttt{GET\ /certificates/\allowbreak{}verify/\{code\}} retorna dados limitados.
Discutir com a equipe se listagem e criação devem ser públicas & Sem
listagem/\allowbreak{}download de certificados de terceiros & \texttt{ADMIN}: acesso
administrativo a certificados; discutir com a equipe se inclui listagem,
PDF e dados pessoais & Listar/\allowbreak{}baixar somente certificados de A & Somente
certificados de B \\
\end{longtable}

\hypertarget{regras-de-campos-e-objetos-a-validar}{%
\subsubsection{Regras de campos e objetos a
validar}\label{regras-de-campos-e-objetos-a-validar}}

\begin{enumerate}
\def\labelenumi{\arabic{enumi}.}
\tightlist
\item
  \texttt{ownerId}, \texttt{userId} e \texttt{role} não podem ser
  escolhidos livremente pelo cliente para mudar a propriedade do
  recurso, assumir identidade alheia ou elevar privilégios. A identidade
  deve vir do token validado; a propriedade deve ser aplicada ou
  conferida no servidor.
\item
  No cadastro/\allowbreak{}edição de curso, uma tentativa de enviar \texttt{ownerId}
  de U2 não deve transferir o curso. Na criação, o proprietário esperado
  é o ator autenticado; na edição, o proprietário não deve ser alterável
  por \texttt{USER}.
\item
  Na criação de atividade, \texttt{sectionId} precisa pertencer a um
  curso que o autor pode editar. Não basta validar que a seção existe.
\item
  \texttt{institutionId}, \texttt{courseId}, \texttt{mediaId}, IDs de
  matrícula, progresso e certificado devem ser revalidados no contexto
  do ator e da operação; UUIDs difíceis de adivinhar não substituem
  autorização.
\item
  A resposta de perfil deve conter apenas campos previstos para o
  próprio ator ou para a função admin autorizada. Campos internos,
  credenciais e hashes nunca devem ser retornados.
\item
  \texttt{correctAnswer} de atividade é campo de autorização/\allowbreak{}leitura:
  não deve ser enviado a anônimos/\allowbreak{}estudantes antes da submissão. O
  tratamento após a submissão deve seguir a regra de feedback da
  atividade.
\item
  Qualquer pessoa, inclusive anônimos e estudantes sem matrícula, pode
  fazer stream de mídia ativa e publicada. Mídias inativas ou não
  publicadas não devem ser transmitidas; metadados e operações de
  escrita continuam restritos ao proprietário ou \texttt{ADMIN}.
\end{enumerate}

\begin{longtable}[]{@{}
  >{\raggedright\arraybackslash}p{(\columnwidth - 2\tabcolsep) * \real{0.5000}}
  >{\raggedright\arraybackslash}p{(\columnwidth - 2\tabcolsep) * \real{0.5000}}@{}}
\toprule\noalign{}
\begin{minipage}[b]{\linewidth}\raggedright
Objeto/\allowbreak{}campo
\end{minipage} & \begin{minipage}[b]{\linewidth}\raggedright
Regra inicial para verificar
\end{minipage} \\
\midrule\noalign{}
\endhead
\bottomrule\noalign{}
\endlastfoot
Usuário: \texttt{role}, \texttt{status} & Somente funções admin
autorizadas; ignorar/\allowbreak{}rejeitar tentativa de alteração por \texttt{/me} ou
endpoint de estudante \\
Curso: \texttt{ownerId} & Derivar do ator na criação e preservar na
edição por \texttt{USER}; não aceitar troca via corpo \\
Atividade: \texttt{sectionId} & \texttt{USER} só pode associar a seção
de curso próprio; \texttt{ADMIN} pode associar conforme regra do
produto \\
Atividade: \texttt{correctAnswer} & Escrita somente por criador/\allowbreak{}admin;
não expor a estudante/\allowbreak{}anônimo antes da submissão \\
Mídia: \texttt{ownerId}, chave de armazenamento e \texttt{status} & Não
controláveis por usuário comum; metadados editáveis pelo
proprietário/\allowbreak{}admin; acesso de leitura separado da edição \\
Matrícula/\allowbreak{}progresso/\allowbreak{}certificado: \texttt{userId} & Derivar do
token/\allowbreak{}registro no servidor; não aceitar ID de outro estudante vindo do
cliente \\
Instituição: nome, domínios e \texttt{isActive} & Escrita somente por
\texttt{ADMIN} \\
\end{longtable}

\hypertarget{contrato-esperado-de-respostas-http}{%
\subsection{Contrato esperado de respostas
HTTP}\label{contrato-esperado-de-respostas-http}}

\begin{longtable}[]{@{}
  >{\raggedright\arraybackslash}p{(\columnwidth - 4\tabcolsep) * \real{0.3333}}
  >{\raggedright\arraybackslash}p{(\columnwidth - 4\tabcolsep) * \real{0.3333}}
  >{\raggedright\arraybackslash}p{(\columnwidth - 4\tabcolsep) * \real{0.3333}}@{}}
\toprule\noalign{}
\begin{minipage}[b]{\linewidth}\raggedright
Situação
\end{minipage} & \begin{minipage}[b]{\linewidth}\raggedright
Resposta esperada
\end{minipage} & \begin{minipage}[b]{\linewidth}\raggedright
Critério de interpretação
\end{minipage} \\
\midrule\noalign{}
\endhead
\bottomrule\noalign{}
\endlastfoot
Endpoint protegido sem token, com token inválido/\allowbreak{}expirado ou identidade
inválida & \textbf{401 Unauthorized} & Não autenticado. Não revelar
dados nem executar a operação. \\
Token válido, mas papel sem permissão para a função & \textbf{403
Forbidden} & Autenticado, porém sem autorização para a função (por
exemplo, \texttt{USER} chamando função exclusiva de \texttt{ADMIN}). \\
Token válido, mas objeto pertence a outro ator & \textbf{403} por
padrão. & O contrato deve ser documentado por endpoint. Se 404 for
escolhido para ocultação, objeto alheio e inexistente devem ser
indistinguíveis na resposta relevante. \\
Recurso realmente inexistente & \textbf{404 Not Found} & Não confundir
ausência do objeto com autorização bem-sucedida. \\
Recurso alheio deliberadamente ocultado & \textbf{404 Not Found} &
Aceitável somente se for política intencional e consistente, sem
conteúdo ou efeito colateral. \\
\end{longtable}

Para todo resultado negado, conferir também o \textbf{estado
antes/\allowbreak{}depois}: nenhum perfil, papel, curso, atividade, mídia, matrícula,
progresso ou certificado pode ter sido alterado. Um código 403/404
isolado não prova que a operação foi bloqueada.

\hypertarget{catuxe1logo-de-casos-de-teste-futuros}{%
\subsection{Catálogo de casos de teste
futuros}\label{catuxe1logo-de-casos-de-teste-futuros}}

Os casos são testes de integração HTTP orientados pela matriz.
\texttt{U1} possui C1/M1; \texttt{U2} possui C2/M2; A e B têm
matrículas, progresso e certificados separados. O status entre
parênteses é alternativa somente quando a política de ocultação de
existência for formalmente escolhida.

\begin{longtable}[]{@{}
  >{\raggedright\arraybackslash}p{(\columnwidth - 4\tabcolsep) * \real{0.3333}}
  >{\raggedright\arraybackslash}p{(\columnwidth - 4\tabcolsep) * \real{0.3333}}
  >{\raggedright\arraybackslash}p{(\columnwidth - 4\tabcolsep) * \real{0.3333}}@{}}
\toprule\noalign{}
\begin{minipage}[b]{\linewidth}\raggedright
ID
\end{minipage} & \begin{minipage}[b]{\linewidth}\raggedright
Ator e cenário / chamada representativa
\end{minipage} & \begin{minipage}[b]{\linewidth}\raggedright
Resultado e verificação de segurança
\end{minipage} \\
\midrule\noalign{}
\endhead
\bottomrule\noalign{}
\endlastfoot
AT-01 & Anônimo: \texttt{GET\ /me/\allowbreak{}profile} e
\texttt{GET\ /student/\allowbreak{}profile} sem \texttt{Authorization} &
\texttt{401}; não retornar perfil. \\
AT-02 & Anônimo: \texttt{POST\ /courses} e \texttt{GET\ /admin/\allowbreak{}users} &
\texttt{401}; nenhum curso/\allowbreak{}consulta administrativa executada. \\
AT-03 & Token ausente, malformado, expirado ou sem \texttt{sub}:
\texttt{GET\ /student/\allowbreak{}enrollments} & \texttt{401}; nenhum dado de
matrícula retornado. \\
AT-04 & \texttt{USER} U1: \texttt{GET\ /admin/\allowbreak{}users} e
\texttt{GET\ /institutions} & \texttt{403}; não retornar
usuários/\allowbreak{}instituições. Repetir com estudante A. \\
AT-05 & \texttt{USER} U1: \texttt{PATCH\ /admin/\allowbreak{}users/\{U1-ou-U2\}/role}
com tentativa de se tornar \texttt{ADMIN} & \texttt{403}; papel
persistido de U1 e de U2 permanece inalterado. \\
AT-06 & \texttt{USER} U1: \texttt{GET\ /courses/\{C2\}} (curso de U2) &
\texttt{403} ; não retornar conteúdo/\allowbreak{}ownerId privado de C2. \\
AT-07 & \texttt{USER} U1: \texttt{PUT\ /courses/\{C2\}} e
\texttt{DELETE\ /courses/\{C2\}} & \texttt{403}; Nega, deve comparar C2
antes/\allowbreak{}depois e comprovar que título, estado e existência não mudaram. \\
AT-08 & \texttt{USER} U1 cria curso enviando no corpo
\texttt{ownerId=U2} e \texttt{role=ADMIN} & A operação autorizada cria
curso com \texttt{ownerId=U1}; campos privilegiados são
ignorados/\allowbreak{}rejeitados e o papel de U1 não muda. \\
AT-09 & \texttt{USER} U1 cria/\allowbreak{}edita atividade em seção de C1 e tenta
criar/\allowbreak{}editar em seção de C2 & Próprio curso: permitir. Outro curso:
negar e não alterar C2. Confirmar pertencimento da seção ao curso no
servidor \texttt{403}. \\
AT-10 & Anônimo: \texttt{GET\ /activities/\{id\}} e
\texttt{GET\ /activities/\allowbreak{}section/\{sectionId\}} para conteúdo não
publicado, além de \texttt{POST\ /activities}; \texttt{USER} U1 não
proprietário: \texttt{GET/\allowbreak{}PUT/\allowbreak{}DELETE\ /activities/\{id\}} de C2 &
Restringir conteúdo privado e escrita: anônimo recebe \texttt{401} na
escrita e não lê atividade não publicada; U1 não lê/\allowbreak{}edita/\allowbreak{}exclui a
atividade de U2.\textbf{Hipótese:} as rotas consultadas não mostram
middleware de autorização; confirmar contrato e resposta diretamente na
API. \\
AT-11 & \texttt{ADMIN}: criar/\allowbreak{}alterar/\allowbreak{}excluir atividade de curso de U2;
comparar com AT-09 & Permite as ações ao \texttt{ADMIN}; verificar que
\texttt{USER} não obtém o mesmo resultado; \textbf{Hipótese:} as rotas
consultadas não mostram checagem de papel/\allowbreak{}propriedade e podem aceitar as
operações de qualquer ator; confirmar via requisições HTTP. Isso não é
uma vulnerabilidade confirmada.. \\
AT-12 & Estudante A: \texttt{PUT\ /student/\allowbreak{}profile} enviando campos
extras \texttt{userId=B} e \texttt{role=ADMIN} & Atualizar, no máximo,
campos permitidos de A; B e os papéis de A/\allowbreak{}B permanecem inalterados. \\
AT-13 & Estudante A: \texttt{GET\ /student/\allowbreak{}enrollments} quando B possui
outra matrícula & \texttt{200} com dados de A apenas; matrícula
exclusiva de B não aparece. Usar curso diferente para detectar mistura
de registros. \\
AT-14 & Estudante A solicita progresso de curso matriculado somente por
B em \texttt{GET\ /student/\allowbreak{}progress/\allowbreak{}courses/\{courseId\}} e tenta
\texttt{userId=B} em query/\allowbreak{}body nos endpoints de progresso &
\texttt{404} conforme contrato aprovado ; em nenhuma resposta aparecem
progresso, pontuação ou conclusão de B. \texttt{userId} deve vir do
token. \\
AT-15 & Estudante A:
\texttt{GET\ /student/\allowbreak{}certificates/\{certificateB\}/pdf} & \texttt{404}
Não encontrará por usar o id do userID passado e do token em uma
comparação. \\
AT-16 & Estudante B lista certificados e progresso; testar também
identificador de A & Listas contêm somente objetos de B. Um ID de A é
negado/\allowbreak{}ocultado, sem campos, PDF ou mudança de estado de A. \\
AT-17 & \texttt{USER} U1: alterar metadados ou excluir M2 em
\texttt{/media/\allowbreak{}images/\{id\}/metadata} e \texttt{/media/\allowbreak{}videos/\{id\}} &
\texttt{403} para objeto alheio; metadados, registro e objeto de
armazenamento permanecem inalterados. Repetir como A contra mídia
privada de U1/U2. \\
AT-18 & Anônimo, A e B: \texttt{GET\ /media/\{id\}/stream} para mídia
ativa/\allowbreak{}publicada e para mídia inativa/\allowbreak{}não publicada & Mídia
ativa/\allowbreak{}publicada: \texttt{200} para todos, sem exigir autenticação ou
matrícula. Mídia inativa/\allowbreak{}não publicada: negar o stream e não revelar o
conteúdo; validar o status HTTP conforme o contrato do endpoint. \\
AT-19 & \texttt{USER} U1:
\texttt{POST/\allowbreak{}PUT/\allowbreak{}DELETE\ /institutions/\{id\}}; \texttt{ADMIN} realiza
operação equivalente com dados de teste & \texttt{USER}: \texttt{403} e
estado inalterado. \texttt{ADMIN}: sucesso somente com payload válido e
dados de teste controlados. \\
AT-20 & \texttt{ADMIN} consulta \texttt{/admin/\allowbreak{}users}, consulta um
usuário e altera o papel de um usuário de teste; em seguida, testar uma
nova requisição desse usuário & Acesso admin permitido. \\
AT-21 & \texttt{USER} U1 consulta \texttt{/me/\allowbreak{}courses} e
\texttt{/me/\allowbreak{}media} usando diferentes valores de paginação, \texttt{kind}
e \texttt{status} & Listagens limitadas a U1; parâmetros de consulta não
removem a condição de propriedade. A tentativa de acessar ID de mídia de
U2 é coberta por AT-17. \\
AT-22 & Anônimo verifica certificado em
\texttt{/certificates/\allowbreak{}verify/\{code\}} com código válido e inválido & A
consulta pública retorna somente atributos explicitamente aprovados para
verificação; não revela e-mail, progresso detalhado ou PDF. Confirmar se
\texttt{userName}, nome do curso, total de seções, data e carga horária
são divulgação intencional. Código inválido retorna \texttt{404}. \\
AT-23 & Anônimo e estudante consultam
\texttt{GET\ /catalog/\allowbreak{}courses/\{C1\}} e inspecionam atividades e campos
retornados & A consulta pública do curso deve retornar 200 com os dados
públicos, mas sem incluir correctAnswer nas atividades.
\textbf{Hipótese}: A rota é publica e retornara o gabarito do curso,
mesmo o usuário nao devendo ter acesso a ele. \\
\end{longtable}

\hypertarget{dados-de-teste-e-controles-positivos}{%
\subsubsection{Dados de teste e controles
positivos}\label{dados-de-teste-e-controles-positivos}}

\begin{itemize}
\tightlist
\item
  Criar U1 e U2 (\texttt{USER}), ADM1 (\texttt{ADMIN}), A e B
  (\texttt{STUDENT}) com credenciais independentes.
\item
  Criar C1 de U1 e C2 de U2, cada um com seção e atividade; configurar
  mídia própria, mídia de curso publicada e mídia inativa.
\item
  Criar matrícula, progresso e certificado para A e B em cursos
  distintos, sem reutilizar IDs.
\item
  Incluir controles positivos: U1 administra C1/M1; U2 administra C2/M2;
  A e B consultam os próprios recursos; ADM1 executa as funções
  administrativas previstas. Eles detectam falsos positivos (negações
  indevidas).
\item
  Para cada tentativa negada, registrar resposta HTTP e consultar
  novamente o banco e os dados de teste por um ator autorizado para
  confirmar ausência de mutação.
\end{itemize}

\hypertarget{muxe9tricas-e-crituxe9rios-de-sucesso}{%
\subsection{Métricas e critérios de
sucesso}\label{muxe9tricas-e-crituxe9rios-de-sucesso}}

\begin{longtable}[]{@{}
  >{\raggedright\arraybackslash}p{(\columnwidth - 4\tabcolsep) * \real{0.3333}}
  >{\raggedright\arraybackslash}p{(\columnwidth - 4\tabcolsep) * \real{0.3333}}
  >{\raggedright\arraybackslash}p{(\columnwidth - 4\tabcolsep) * \real{0.3333}}@{}}
\toprule\noalign{}
\begin{minipage}[b]{\linewidth}\raggedright
Métrica
\end{minipage} & \begin{minipage}[b]{\linewidth}\raggedright
Fórmula
\end{minipage} & \begin{minipage}[b]{\linewidth}\raggedright
Critério de sucesso proposto
\end{minipage} \\
\midrule\noalign{}
\endhead
\bottomrule\noalign{}
\endlastfoot
Cobertura de regras & Regras com ao menos um teste positivo e um
negativo ÷ regras aplicáveis inventariadas × 100 & \textbf{100\%} \\
Cobertura de funções & Pares método HTTP--endpoint sensíveis com
autorização testada ÷ pares sensíveis inventariados × 100 &
\textbf{100\%} \\
Cobertura de objetos & Combinações próprias, alheias e inexistentes
previstas e exercitadas ÷ combinações planejadas aplicáveis × 100 &
\textbf{100\%} \\
Cobertura de campos & Campos sensíveis inventariados com leitura/\allowbreak{}escrita
permitida e negada testadas ÷ campos sensíveis inventariados × 100 &
\textbf{100\%} \\
Tentativas negadas bloqueadas & Tentativas não autorizadas sem
dado/\allowbreak{}efeito indevido ÷ total de tentativas não autorizadas executadas ×
100 & \textbf{100\%} \\
Bypasses confirmados & Número de operações/\allowbreak{}dados não autorizados que
foram obtidos ou alterados & \textbf{0} \\
Falsos positivos & Operações legítimas de controle negadas ÷ operações
legítimas de controle executadas × 100 & \textbf{0\%} \\
Integridade após negação & Tentativas negadas sem mutação persistida ÷
tentativas negadas com estado verificado × 100 & \textbf{100\%} \\
Rastreabilidade ASVS & Requisitos ASVS aplicáveis com ao menos um caso
associado ÷ requisitos ASVS aplicáveis selecionados × 100 &
\textbf{100\%} \\
Severidade & Contagem de falhas confirmadas por severidade após
validação e análise de impacto & Registrar \textbf{100\%} das falhas com
severidade; \textbf{0 críticas/\allowbreak{}altas abertas} como condição para aceite
futuro \\
\end{longtable}

\hypertarget{severidade-proposta-para-triagem}{%
\subsubsection{Severidade proposta para
triagem}\label{severidade-proposta-para-triagem}}

\begin{itemize}
\tightlist
\item
  \textbf{Crítica:} acesso administrativo/\allowbreak{}sistêmico sem autorização ou
  alteração ampla de dados/\allowbreak{}contas.
\item
  \textbf{Alta:} leitura ou alteração não autorizada de dados pessoais,
  certificados, progresso, papéis ou conteúdo de múltiplos usuários.
\item
  \textbf{Média:} acesso/\allowbreak{}alteração indevida de escopo limitado, sem
  evidência de impacto amplo.
\item
  \textbf{Baixa:} exposição de informação de baixo impacto ou falha com
  efeito muito restrito.
\end{itemize}

Classificar somente após reproduzir o resultado no ambiente autorizado,
registrar evidência e estimar impacto. Uma hipótese derivada de leitura
de código não recebe severidade como se fosse achado confirmado.

\hypertarget{rastreabilidade-ao-owasp-asvs}{%
\subsection{Rastreabilidade ao OWASP
ASVS}\label{rastreabilidade-ao-owasp-asvs}}

Referência escolhida: \textbf{OWASP ASVS 5.0.0}, capítulo V8
Authorization.

\begin{longtable}[]{@{}
  >{\raggedright\arraybackslash}p{(\columnwidth - 4\tabcolsep) * \real{0.3333}}
  >{\raggedright\arraybackslash}p{(\columnwidth - 4\tabcolsep) * \real{0.3333}}
  >{\raggedright\arraybackslash}p{(\columnwidth - 4\tabcolsep) * \real{0.3333}}@{}}
\toprule\noalign{}
\begin{minipage}[b]{\linewidth}\raggedright
ASVS 5.0.0
\end{minipage} & \begin{minipage}[b]{\linewidth}\raggedright
Aplicação ao planejamento
\end{minipage} & \begin{minipage}[b]{\linewidth}\raggedright
Casos relacionados
\end{minipage} \\
\midrule\noalign{}
\endhead
\bottomrule\noalign{}
\endlastfoot
V8.1.1 & Documentar regras de autorização por função e dados/\allowbreak{}objetos &
Matriz; AT-04 a AT-11, AT-19 a AT-22 \\
V8.1.2 & Documentar restrições de leitura e escrita por campo & Regras
de campos; AT-05, AT-08, AT-12, AT-14, AT-22, AT-23 \\
V8.2.1 & Restringir funções a consumidores com permissão explícita &
AT-01 a AT-05, AT-10, AT-11, AT-19, AT-20 \\
V8.2.2 & Restringir dados específicos e mitigar IDOR/\allowbreak{}BOLA & AT-06,
AT-07, AT-09, AT-13 a AT-18, AT-21 \\
V8.2.3 & Restringir acesso a propriedades/\allowbreak{}campos e mitigar BOPLA &
AT-05, AT-08, AT-12, AT-14, AT-22, AT-23 \\
V8.3.1 & Aplicar autorização em camada confiável do servidor, não só no
cliente & Todos; chamadas diretas à \texttt{educado-api} \\
V8.3.2 & Aplicar mudanças de permissão de forma efetiva, inclusive em
sessões/\allowbreak{}tokens & AT-05 e AT-20; política de revogação/\allowbreak{}expiração a
confirmar \\
V8.4.1 & Controles entre instituições/\allowbreak{}tenants & AT-19 e casos adicionais
após confirmar o modelo multi-instituição \\
\end{longtable}

V8.1.3, V8.1.4, V8.2.4 e V8.4.2 tratam atributos contextuais/\allowbreak{}adaptativos
ou controles administrativos de nível mais amplo e ficam fora do escopo
inicial, salvo requisito do produto que os torne aplicáveis. A cobertura
não deve ser declarada para requisitos excluídos.

\hypertarget{hipuxf3teses-do-cuxf3digo-a-confirmar-nuxe3o-suxe3o-vulnerabilidades-confirmadas}{%
\subsection{Hipóteses do código a confirmar (não são vulnerabilidades
confirmadas)}\label{hipuxf3teses-do-cuxf3digo-a-confirmar-nuxe3o-suxe3o-vulnerabilidades-confirmadas}}

Na leitura estática usada para definir os cenários, foram observados
pontos que justificam testes, mas não provam falha em execução:

\begin{itemize}
\tightlist
\item
  Há verificações de papel no middleware e nas funções administrativas;
  há também checagens de propriedade em serviços de curso e de mídia. Os
  testes devem verificar as rotas e os efeitos completos, não apenas a
  existência desses helpers.
\item
  As rotas REST de atividades consultadas não mostram middleware de
  autorização no próprio arquivo. AT-09 a AT-11 validarão a política
  esperada de dono do curso/\allowbreak{}admin e confirmarão como os endpoints se
  comportam via API.
\item
  A projeção de atividades no detalhe público do catálogo inclui o
  atributo \texttt{correctAnswer}, e a verificação pública do
  certificado retorna \texttt{userName}, curso, total de seções, data e
  carga horária. Isso motiva AT-22/AT-23 para confirmar a política de
  divulgação campo a campo; não é uma conclusão de vulnerabilidade.
\item
  O stream de mídia tem lógica própria de token e trata estudantes
  diferentemente de leitores de metadados. AT-18 verificará que qualquer
  pessoa consegue transmitir mídia ativa e publicada, independentemente
  de autenticação ou matrícula, e que mídia inativa/\allowbreak{}não publicada não é
  transmitida.
\item
  O mapeamento atual distingue \texttt{USER}, \texttt{ADMIN} e
  \texttt{STUDENT}. Para esta proposta, \texttt{USER} significa autor de
  conteúdo e estudantes A/\allowbreak{}B são contas diferentes com papel
  \texttt{STUDENT}.
\end{itemize}

Qualquer desvio deve ser reproduzido, comparado ao contrato aprovado,
confirmado quanto à ausência/\allowbreak{}presença de efeito e registrado como
resultado somente depois da execução autorizada.

\hypertarget{pruxe9-requisitos-limites-e-seguranuxe7a-operacional}{%
\subsection{Pré-requisitos, limites e segurança
operacional}\label{pruxe9-requisitos-limites-e-seguranuxe7a-operacional}}

\begin{itemize}
\tightlist
\item
  Executar \textbf{somente em ambiente local ou homologação com
  autorização explícita}; nunca contra produção, usuários reais ou
  serviços de terceiros.
\item
  Usar banco e armazenamento de mídia isolados, contas sintéticas e IDs
  gerados para teste. Não reutilizar tokens, segredos ou dados pessoais
  reais.
\item
  Preparar backups e limpeza dos dados de teste, além de registrar os
  estados iniciais para provar que tentativas negadas não alteraram
  dados.
\item
  Manter alcance limitado aos endpoints e cenários acordados. Não
  realizar força bruta, varredura indiscriminada ou tentativas contra
  infraestrutura externa.
\item
  Confirmar previamente a ocultação de objetos com 404, administração de
  matrículas/\allowbreak{}progresso/\allowbreak{}certificados, escopo de acesso de ADMIN e
  isolamento por instituição.
\item
  Este plano não implementa nem executa os testes e não apresenta
  resultados experimentais.
\end{itemize}

\hypertarget{fontes-consultadas-nesta-seuxe7uxe3o}{%
\subsection{Fontes consultadas nesta
seção}\label{fontes-consultadas-nesta-seuxe7uxe3o}}

\begin{itemize}
\tightlist
\item
  OWASP. \textbf{Application Security Verification Standard (ASVS),
  versão estável 5.0.0}, V8 Authorization:
  \href{https://github.com/OWASP/ASVS/blob/v5.0.0_release/5.0/en/0x17-V8-Authorization.md}{requisitos
  V8} e
  \href{https://github.com/OWASP/ASVS/releases/tag/v5.0.0_release}{release
  estável}.
\item
  OWASP. \textbf{Authorization Cheat Sheet}:
  \href{https://cheatsheetseries.owasp.org/cheatsheets/Authorization_Cheat_Sheet.html}{cheatsheetseries.owasp.org}.
\item
  Código consultado para delimitar as superfícies (sem concluir
  vulnerabilidade):
  \href{../educado-api/src/interface/http/middlewares/auth-jwt.ts}{middleware
  de autenticação},
  \href{../educado-api/src/interface/http/middlewares/require-role.ts}{middleware
  de papel},
  \href{../educado-api/src/application/courses/course-access-service.ts}{acesso
  a cursos},
  \href{../educado-api/src/application/media/media-access-service.ts}{acesso
  a mídias},
  \href{../educado-api/src/routes/activities/activities.ts}{rotas de
  atividades}, \href{../educado-api/src/routes/student/student.ts}{rotas
  de estudante} e
  \href{../educado-api/src/routes/admin/registrations.ts}{rotas
  administrativas}.
\end{itemize}

\hypertarget{limitauxe7uxf5es-e-conclusuxe3o}{%
\section{Limitações e conclusão}\label{limitauxe7uxf5es-e-conclusuxe3o}}

\hypertarget{limitauxe7uxf5es}{%
\subsection{Limitações}\label{limitauxe7uxf5es}}

O trabalho se limita à pesquisa e ao planejamento. As técnicas,
ferramentas e métricas descritas não foram aplicadas ao Educado nesta
entrega, então nenhuma afirmação deste documento deve ser lida como
resultado de teste. A matriz de autorização proposta parte da leitura do
código e da documentação do projeto e precisa ser validada contra o
comportamento real da \texttt{educado-api}, que é onde a decisão de
acesso efetivamente acontece. A execução futura deve ocorrer apenas em
ambiente local ou de homologação autorizado, com dados descartáveis.

\hypertarget{conclusuxe3o}{%
\subsection{Conclusão}\label{conclusuxe3o}}

Autorização pode ser testada de forma sistemática quando a política é
explicitada como uma matriz de sujeito, recurso, operação e resultado
esperado. As referências estudadas se complementam: o OWASP Top 10
indica a relevância do risco, o ASVS fornece requisitos verificáveis, o
WSTG descreve procedimentos de teste e o CWE classifica as fraquezas
encontradas {[}2{]}. Para o Educado, a equipe propõe combinar
Playwright, para cenários conhecidos e repetíveis entre perfis, com o
OWASP ZAP, para a varredura de rotas, usando Vitest apenas como
complemento no frontend. A próxima etapa do seminário é executar o plano
em ambiente controlado e medir os resultados com as métricas definidas.

\hypertarget{bibliography}{%
\section*{Referências}\label{bibliography}}
\addcontentsline{toc}{section}{Referências}

\hypertarget{refs}{}
\begin{CSLReferences}{0}{0}
\leavevmode\vadjust pre{\hypertarget{ref-educado-repos}{}}%
\CSLLeftMargin{{[}1{]} }%
\CSLRightInline{ErasmusEgalitarian, {``Repositórios do projeto Educado:
educado-api, educado-web, educado-app e educado-docs''}. GitHub, 2026.
Acessado: 12 de novembro de 2026. {[}Online{]}. Disponível em:
\url{https://github.com/ErasmusEgalitarian}}

\leavevmode\vadjust pre{\hypertarget{ref-owasp-top10-2025}{}}%
\CSLLeftMargin{{[}2{]} }%
\CSLRightInline{OWASP Foundation, {``A01:2025 Broken Access Control''}.
Acessado: 12 de novembro de 2026. {[}Online{]}. Disponível em:
\url{https://top10.owasp.org/2025/A01_2025-Broken_Access_Control/index.html}}

\leavevmode\vadjust pre{\hypertarget{ref-nist-sp800-162}{}}%
\CSLLeftMargin{{[}3{]} }%
\CSLRightInline{V. C. Hu \emph{et al.}, {``Guide to Attribute Based
Access Control ({ABAC}) Definition and Considerations''}, National
Institute of Standards; Technology, NIST Special Publication 800-162,
2014. doi:
\href{https://doi.org/10.6028/NIST.SP.800-162}{10.6028/NIST.SP.800-162}.}

\leavevmode\vadjust pre{\hypertarget{ref-saltzer-schroeder-1975}{}}%
\CSLLeftMargin{{[}4{]} }%
\CSLRightInline{J. H. Saltzer e M. D. Schroeder, {``The Protection of
Information in Computer Systems''}, \emph{Proceedings of the IEEE}, vol.
63, nº 9, p. 1278--1308, 1975, doi:
\href{https://doi.org/10.1109/PROC.1975.9939}{10.1109/PROC.1975.9939}.}

\leavevmode\vadjust pre{\hypertarget{ref-owasp-authorization}{}}%
\CSLLeftMargin{{[}5{]} }%
\CSLRightInline{OWASP Foundation, {``Authorization Cheat Sheet''}.
Acessado: 12 de novembro de 2026. {[}Online{]}. Disponível em:
\url{https://cheatsheetseries.owasp.org/cheatsheets/Authorization_Cheat_Sheet.html}}

\leavevmode\vadjust pre{\hypertarget{ref-ferraiolo-kuhn-1992}{}}%
\CSLLeftMargin{{[}6{]} }%
\CSLRightInline{D. F. Ferraiolo e D. R. Kuhn, {``Role-Based Access
Controls''}, em \emph{Proceedings of the 15th National Computer Security
Conference}, Baltimore, MD, 1992, p. 554--563. Acessado: 12 de novembro
de 2026. {[}Online{]}. Disponível em:
\url{https://csrc.nist.gov/CSRC/media/Publications/conference-paper/1992/10/13/role-based-access-controls/documents/ferraiolo-kuhn-92.pdf}}

\leavevmode\vadjust pre{\hypertarget{ref-rel02-iso25010-2023}{}}%
\CSLLeftMargin{{[}7{]} }%
\CSLRightInline{ISO/\allowbreak{}IEC, {``Systems and software engineering --- Systems
and software Quality Requirements and Evaluation ({SQuaRE}) --- Product
quality model''}, International Organization for Standardization,
ISO/\allowbreak{}IEC 25010:2023, 2023. Acessado: 12 de novembro de 2026.
{[}Online{]}. Disponível em:
\url{https://www.iso.org/standard/78176.html}}

\leavevmode\vadjust pre{\hypertarget{ref-rel02-top10-2025}{}}%
\CSLLeftMargin{{[}8{]} }%
\CSLRightInline{OWASP Foundation, {``{OWASP Top 10:2025}''}. 2025.
Acessado: 12 de novembro de 2026. {[}Online{]}. Disponível em:
\url{https://top10.owasp.org/2025/}}

\leavevmode\vadjust pre{\hypertarget{ref-rel02-top10-a01-2025}{}}%
\CSLLeftMargin{{[}9{]} }%
\CSLRightInline{OWASP Foundation, {``{OWASP Top 10:2025} --- {A01:2025
Broken Access Control}''}. 2025. Acessado: 12 de novembro de 2026.
{[}Online{]}. Disponível em:
\url{https://top10.owasp.org/2025/A01_2025-Broken_Access_Control/index.html}}

\leavevmode\vadjust pre{\hypertarget{ref-rel02-asvs-5}{}}%
\CSLLeftMargin{{[}10{]} }%
\CSLRightInline{OWASP Foundation, {``{Application Security Verification
Standard 5.0.0} --- {V8 Authorization}''}. 2025. Acessado: 12 de
novembro de 2026. {[}Online{]}. Disponível em:
\url{https://github.com/OWASP/ASVS/blob/v5.0.0/5.0/en/0x17-V8-Authorization.md}}

\leavevmode\vadjust pre{\hypertarget{ref-rel02-wstg42-athz01}{}}%
\CSLLeftMargin{{[}11{]} }%
\CSLRightInline{OWASP Foundation, {``{Web Security Testing Guide 4.2}
--- Testing Directory Traversal File Include''}. 2020. Acessado: 12 de
novembro de 2026. {[}Online{]}. Disponível em:
\url{https://wstg.owasp.org/v4.2/4-Web_Application_Security_Testing/05-Authorization_Testing/01-Testing_Directory_Traversal_File_Include/}}

\leavevmode\vadjust pre{\hypertarget{ref-rel02-wstg42-athz02}{}}%
\CSLLeftMargin{{[}12{]} }%
\CSLRightInline{OWASP Foundation, {``{Web Security Testing Guide 4.2}
--- Testing for Bypassing Authorization Schema''}. 2020. Acessado: 12 de
novembro de 2026. {[}Online{]}. Disponível em:
\url{https://wstg.owasp.org/v4.2/4-Web_Application_Security_Testing/05-Authorization_Testing/02-Testing_for_Bypassing_Authorization_Schema/}}

\leavevmode\vadjust pre{\hypertarget{ref-rel02-wstg42-athz03}{}}%
\CSLLeftMargin{{[}13{]} }%
\CSLRightInline{OWASP Foundation, {``{Web Security Testing Guide 4.2}
--- Testing for Privilege Escalation''}. 2020. Acessado: 12 de novembro
de 2026. {[}Online{]}. Disponível em:
\url{https://wstg.owasp.org/v4.2/4-Web_Application_Security_Testing/05-Authorization_Testing/03-Testing_for_Privilege_Escalation/}}

\leavevmode\vadjust pre{\hypertarget{ref-rel02-wstg42-athz04}{}}%
\CSLLeftMargin{{[}14{]} }%
\CSLRightInline{OWASP Foundation, {``{Web Security Testing Guide 4.2}
--- Testing for Insecure Direct Object References''}. 2020. Acessado: 12
de novembro de 2026. {[}Online{]}. Disponível em:
\url{https://wstg.owasp.org/v4.2/4-Web_Application_Security_Testing/05-Authorization_Testing/04-Testing_for_Insecure_Direct_Object_References/}}

\leavevmode\vadjust pre{\hypertarget{ref-rel02-api1-2023}{}}%
\CSLLeftMargin{{[}15{]} }%
\CSLRightInline{OWASP Foundation, {``{OWASP API Security Top 10:2023}
--- {API1:2023 Broken Object Level Authorization}''}. 2023. Acessado: 12
de novembro de 2026. {[}Online{]}. Disponível em:
\url{https://api-security.owasp.org/editions/2023/en/0xa1-broken-object-level-authorization/}}

\leavevmode\vadjust pre{\hypertarget{ref-rel02-api3-2023}{}}%
\CSLLeftMargin{{[}16{]} }%
\CSLRightInline{OWASP Foundation, {``{OWASP API Security Top 10:2023}
--- {API3:2023 Broken Object Property Level Authorization}''}. 2023.
Acessado: 12 de novembro de 2026. {[}Online{]}. Disponível em:
\url{https://api-security.owasp.org/editions/2023/en/0xa3-broken-object-property-level-authorization/}}

\leavevmode\vadjust pre{\hypertarget{ref-rel02-api5-2023}{}}%
\CSLLeftMargin{{[}17{]} }%
\CSLRightInline{OWASP Foundation, {``{OWASP API Security Top 10:2023}
--- {API5:2023 Broken Function Level Authorization}''}. 2023. Acessado:
12 de novembro de 2026. {[}Online{]}. Disponível em:
\url{https://api-security.owasp.org/editions/2023/en/0xa5-broken-function-level-authorization/}}

\leavevmode\vadjust pre{\hypertarget{ref-rel02-cwe5}{}}%
\CSLLeftMargin{{[}18{]} }%
\CSLRightInline{MITRE, {``{Common Weakness Enumeration (CWE)}, Version
5.0''}. 2026. Acessado: 12 de novembro de 2026. {[}Online{]}. Disponível
em: \url{https://cwe.mitre.org/data/xml/cwec_v5.0.xml.zip}}

\leavevmode\vadjust pre{\hypertarget{ref-rel02-cwe862}{}}%
\CSLLeftMargin{{[}19{]} }%
\CSLRightInline{MITRE, {``{CWE-862: Missing Authorization}''}. 2026.
Acessado: 12 de novembro de 2026. {[}Online{]}. Disponível em:
\url{https://cwe.mitre.org/data/definitions/862.html}}

\leavevmode\vadjust pre{\hypertarget{ref-rel02-cwe863}{}}%
\CSLLeftMargin{{[}20{]} }%
\CSLRightInline{MITRE, {``{CWE-863: Incorrect Authorization}''}. 2026.
Acessado: 12 de novembro de 2026. {[}Online{]}. Disponível em:
\url{https://cwe.mitre.org/data/definitions/863.html}}

\leavevmode\vadjust pre{\hypertarget{ref-rel02-cwe639}{}}%
\CSLLeftMargin{{[}21{]} }%
\CSLRightInline{MITRE, {``{CWE-639: Authorization Bypass Through
User-Controlled Key}''}. 2026. Acessado: 12 de novembro de 2026.
{[}Online{]}. Disponível em:
\url{https://cwe.mitre.org/data/definitions/639.html}}

\leavevmode\vadjust pre{\hypertarget{ref-owasp-wstg}{}}%
\CSLLeftMargin{{[}22{]} }%
\CSLRightInline{OWASP Foundation, {``Web Security Testing Guide:
Authorization Testing''}. Acessado: 12 de novembro de 2026.
{[}Online{]}. Disponível em:
\url{https://owasp.org/www-project-web-security-testing-guide/latest/4-Web_Application_Security_Testing/05-Authorization_Testing/}}

\leavevmode\vadjust pre{\hypertarget{ref-owasp-api1}{}}%
\CSLLeftMargin{{[}23{]} }%
\CSLRightInline{OWASP Foundation, {``API1:2023 Broken Object Level
Authorization''}. Acessado: 12 de novembro de 2026. {[}Online{]}.
Disponível em:
\url{https://owasp.org/API-Security/editions/2023/en/0xa1-broken-object-level-authorization/}}

\leavevmode\vadjust pre{\hypertarget{ref-owasp-api3}{}}%
\CSLLeftMargin{{[}24{]} }%
\CSLRightInline{OWASP Foundation, {``API3:2023 Broken Object Property
Level Authorization''}. Acessado: 12 de novembro de 2026. {[}Online{]}.
Disponível em:
\url{https://owasp.org/API-Security/editions/2023/en/0xa3-broken-object-property-level-authorization/}}

\leavevmode\vadjust pre{\hypertarget{ref-openapi-spec}{}}%
\CSLLeftMargin{{[}25{]} }%
\CSLRightInline{OpenAPI Initiative, {``OpenAPI Specification''}.
Acessado: 12 de novembro de 2026. {[}Online{]}. Disponível em:
\url{https://spec.openapis.org/oas/latest.html}}

\leavevmode\vadjust pre{\hypertarget{ref-zap}{}}%
\CSLLeftMargin{{[}26{]} }%
\CSLRightInline{ZAP by Checkmarx, {``OWASP ZAP''}. Acessado: 12 de
novembro de 2026. {[}Online{]}. Disponível em:
\url{https://www.zaproxy.org/}}

\leavevmode\vadjust pre{\hypertarget{ref-zap-access-control}{}}%
\CSLLeftMargin{{[}27{]} }%
\CSLRightInline{ZAP by Checkmarx, {``Access Control Testing''}.
Acessado: 12 de novembro de 2026. {[}Online{]}. Disponível em:
\url{https://www.zaproxy.org/docs/desktop/addons/access-control-testing/}}

\leavevmode\vadjust pre{\hypertarget{ref-playwright}{}}%
\CSLLeftMargin{{[}28{]} }%
\CSLRightInline{Microsoft, {``Playwright documentation''}. Acessado: 12
de novembro de 2026. {[}Online{]}. Disponível em:
\url{https://playwright.dev/docs/intro}}

\leavevmode\vadjust pre{\hypertarget{ref-playwright-api-request}{}}%
\CSLLeftMargin{{[}29{]} }%
\CSLRightInline{Microsoft, {``{APIRequestContext}''}. Acessado: 12 de
novembro de 2026. {[}Online{]}. Disponível em:
\url{https://playwright.dev/docs/api/class-apirequestcontext}}

\leavevmode\vadjust pre{\hypertarget{ref-vitest}{}}%
\CSLLeftMargin{{[}30{]} }%
\CSLRightInline{Vitest contributors, {``Vitest: Getting Started''}.
Acessado: 12 de novembro de 2026. {[}Online{]}. Disponível em:
\url{https://vitest.dev/guide/}}

\leavevmode\vadjust pre{\hypertarget{ref-schemathesis}{}}%
\CSLLeftMargin{{[}31{]} }%
\CSLRightInline{Schemathesis contributors, {``Schemathesis
documentation''}. Acessado: 12 de novembro de 2026. {[}Online{]}.
Disponível em: \url{https://schemathesis.readthedocs.io/en/stable/}}

\leavevmode\vadjust pre{\hypertarget{ref-burp-community}{}}%
\CSLLeftMargin{{[}32{]} }%
\CSLRightInline{PortSwigger, {``Burp Suite Community Edition''}.
Acessado: 12 de novembro de 2026. {[}Online{]}. Disponível em:
\url{https://portswigger.net/burp/communitydownload}}

\leavevmode\vadjust pre{\hypertarget{ref-autorize}{}}%
\CSLLeftMargin{{[}33{]} }%
\CSLRightInline{B. Tawily, {``Autorize: automatic authorization
enforcement detection extension for Burp Suite''}. Acessado: 12 de
novembro de 2026. {[}Online{]}. Disponível em:
\url{https://github.com/PortSwigger/autorize}}

\end{CSLReferences}
